\documentclass[10pt,a4]{examenEPN0}
\usepackage{comandosEPN}

\usepackage{empheq}

\newcommand{\itemcp}{\stepcounter{enumi}\item[ \textcolor{red}{\theenumi.$^*$}] }

\newcommand{\iitemcp}{\stepcounter{enumii}\item[ \textcolor{red}{\theenumii)$^*$}] }


\materia{Álgebra Lineal}
%\autor{Departamento de Formación Básica}
\autor{GR \hspace{1cm}.}
%\fecha{Semestre 2026-A}
\fecha{Nombre: \hspace{9cm} Cédula:}
\logouno[1.6cm]{Logo_EPN.eps}
\logodos[2.6cm]{Logo_DFB.eps}

\newcommand{\la}{\lambda}

%Datos
\newcommand{\datosAlgebra}[2]{
\EnExamen
{
\vspace{5mm}
    \noindent
    Nombres: \hrulefill\
    Carrera: \rule{4.5cm}{0.4pt}\
    \ifthenelse{\equal{#1}{}}%
    {Paralelo: \underline{\parbox{1.8cm}{\parbox{1cm}{}}}\\[3mm]}
    {Paralelo: \underline{\parbox{1.8cm}{\centering #1}}\\[3mm]}
    C\'edula: \rule{5cm}{0.4pt}\
    Firma: \hrulefill\
    Fecha: \underline{\parbox{2.1cm}{\centering #2}}
}}


% \newcommand{\indicacionesPrueba}{%
% \EnExamen{
% \begin{indicaciones}
%     \item La prueba tiene una duraci\'on  de \textbf{120 minutos.}
%     \item No se permite la utilizaci\'on de dispositivos electr\'onicos (celular, tablets, calculadora, etc.)
%     \item Leer detenidamente los enunciados, \textbf{justificar la respuesta para acceder al puntaje total}.
%     \item La prueba es completamente individual. Los estudiantes deberán dejar sus teléfonos celulares guardados y no podrán tenerlos consigo (por ejemplo, en los bolsillos de su ropa) durante toda la evaluación. En caso de detectarse algún acto de deshonestidad académica o el incumplimiento de esta norma, no podrán seguir rindiendo la prueba y su calificación será \textbf{cero}.
% \end{indicaciones}
% \pagestyle{empty}
% }}

\newcommand{\indicacionesPrueba}{%
\EnExamen{
\begin{indicaciones}
    \item La prueba tiene una duraci\'on  de \textbf{60 minutos.}
    \item Leer detenidamente los enunciados, \textbf{justificar la respuesta} para acceder al puntaje total.
    \item No se permite la utilizaci\'on de dispositivos electr\'onicos (celular, tablets, calculadora, etc.)
    \item La prueba es completamente individual
    \item No se permite portar teléfonos celulares, relojes inteligentes ni otros dispositivos electrónicos durante el examen, aun cuando permanezcan en los bolsillos. Estos dispositivos deberán guardarse en la mochila antes del inicio de la evaluación y esta deberá mantenerse cerrada durante todo el examen.
    \item He leído las indicaciones y acepto que el incumplimiento de estas se penaliza con \textbf{cero}. Firma:...................................

    
\end{indicaciones}
\pagestyle{empty}
}}

\newcommand{\indicacionesExamen}{
\EnExamen{\begin{indicaciones}
    \item La evaluación tiene una duraci\'on  de \textbf{120 minutos.}
    \item No se permite la utilizaci\'on de dispositivos electr\'onicos (celular, tablets, calculadora, etc.)
    %\item Leer detenidamente los enunciados, justificar la respuesta para acceder al puntaje total.
    \item El examen es completamente individual. En caso de detectarse alg\'un acto de deshonestidad acad\'emica no podr\'a seguir rindiendo la prueba y su calificaci\'on ser\'a \textbf{cero}.
\end{indicaciones}
}\pagestyle{empty}
}

%% Imprimir las una hoja A4, dos veces en A5
\newcommand{\imprimirDoble}
{%
\usepackage{pgfmorepages}
\pgfmorepagesloadextralayouts
\pgfpagesuselayout{repeated 2-up}[a4paper,landscape]
\thispagestyle{empty}
}


%Este sí se puede comentar
\newenvironment{solucion}{\color{blue}\proof[Solución]}{\endproof\color{black}}




% \renewcommand{\@autor}{
% ---\\
% \phantom{.} 
% ...}
% \makeatother

\examen{Prueba 02: solucionario}
\makeatletter
\renewcommand{\@autor}{Departamento de Formación Básica}
\renewcommand{\@fecha}{GR6M -- Viernes de 11 a 13 -- 06/11/2026}
\makeatother
\begin{document}
\encabezado
\noindent\textbf{Criterio general.} Se asigna crédito parcial por los pasos
correctos. Un error aritmético se penaliza donde se origina; los pasos
posteriores se valoran por su coherencia, sin repetir la penalización.
Se aceptan procedimientos equivalentes que respeten el método solicitado.

\section*{1. Sistema homogéneo singular (2 puntos)}
\noindent\textbf{Rúbrica.}
\begin{itemize}[itemsep=2pt]
    \item Reconocer que el sistema homogéneo es compatible: 0,5 puntos.
    \item Deducir de $\det(A)=0$ que $\rang(A)<4$: 0,5 puntos.
    \item Justificar que hay variables libres e infinitas soluciones: 0,5 puntos.
    \item Seleccionar la opción (c): 0,5 puntos.
\end{itemize}
El sistema $Ax=0$ siempre admite la solución trivial $x=0$, por lo que
es compatible. Además, $\det(A)=0$ implica $\rang(A)<4$. Como la columna
añadida es nula, $\rang(A\mid0)=\rang(A)<4$. Por el teorema de
Rouché--Frobenius, quedan variables libres y existen infinitas soluciones
reales. La respuesta correcta es la \textbf{opción (c)}.

\section*{2. Procesos del centro de datos (6 puntos)}
\noindent\textbf{Rúbrica.}
\begin{itemize}[itemsep=2pt]
    \item Definir las tres variables como cantidades de procesos: 0,5 puntos.
    \item Plantear las ecuaciones de recursos: 1,5 puntos (0,5 por ecuación).
    \item Identificar $A$, $x$ y $b$: 1 punto.
    \item Calcular $\det(A)$ y comprobar que Cramer es aplicable: 0,5 puntos.
    \item Construir las tres matrices $A^{(i)}$: 0,75 puntos (0,25 por matriz).
    \item Calcular los tres determinantes de reemplazo: 0,75 puntos (0,25 por determinante).
    \item Aplicar los tres cocientes de Cramer: 0,75 puntos (0,25 por incógnita).
    \item Interpretar las cantidades obtenidas: 0,25 puntos.
\end{itemize}
\noindent\textbf{a) Modelo (3 puntos).}
Sean $x_A$, $x_B$ y $x_C$ las cantidades de procesos de análisis,
respaldo e inteligencia artificial, respectivamente. Deben ser enteros no negativos.
El enunciado no asigna consumo de almacenamiento al proceso B; en el modelo
su coeficiente es 0. Las ecuaciones son
\[
    \left\{\begin{aligned}
        2x_A+3x_B+2x_C&=21 &&\text{(CPU)},\\
        x_A+2x_B+x_C&=12 &&\text{(RAM)},\\
        x_A+2x_C&=10 &&\text{(almacenamiento)}.
    \end{aligned}\right.
\]
Así, $Ax=b$, con
\[
    A=\begin{pmatrix}2&3&2\\1&2&1\\1&0&2\end{pmatrix},\quad
    x=\begin{pmatrix}x_A\\x_B\\x_C\end{pmatrix},\quad
    b=\begin{pmatrix}21\\12\\10\end{pmatrix}.
\]
\noindent\textbf{b) Regla de Cramer (3 puntos).}
Por expansión en la primera fila,
\[
    \det(A)=2(4-0)-3(2-1)+2(0-2)=8-3-4=1\ne0.
\]
Definimos $A^{(i)}$ reemplazando la columna $i$ de $A$ por $b$:
\[
    A^{(1)}=\begin{pmatrix}21&3&2\\12&2&1\\10&0&2\end{pmatrix},\quad
    A^{(2)}=\begin{pmatrix}2&21&2\\1&12&1\\1&10&2\end{pmatrix},\quad
    A^{(3)}=\begin{pmatrix}2&3&21\\1&2&12\\1&0&10\end{pmatrix}.
\]
\begin{align*}
    \det(A^{(1)})&=21(4)-3(24-10)+2(0-20)=84-42-40=2,\\
    \det(A^{(2)})&=2(24-10)-21(2-1)+2(10-12)=28-21-4=3,\\
    \det(A^{(3)})&=2(20)-3(10-12)+21(0-2)=40+6-42=4.
\end{align*}
Por tanto,
\[
    x_A=\frac{\det(A^{(1)})}{\det(A)}=2,\qquad
    x_B=\frac{\det(A^{(2)})}{\det(A)}=3,\qquad
    x_C=\frac{\det(A^{(3)})}{\det(A)}=4.
\]
\textbf{Se ejecutaron 2 procesos de análisis, 3 de respaldo y 4 de inteligencia artificial.}
En efecto, $4+9+8=21$, $2+6+4=12$ y $2+8=10$.

\newpage
\section*{3. Por qué $B$ no es la inversa de $A$ (6 puntos)}
\noindent\textbf{Rúbrica.}
\begin{itemize}[itemsep=2pt]
    \item Indicar que ser inversa exige $AB=I_2$: 1 punto.
    \item Calcular correctamente la segunda columna de $AB$: 2 puntos
          (1 por entrada).
    \item Comparar con la segunda columna de $I_2$ y plantear las condiciones: 1 punto.
    \item Deducir $\alpha=2$, $\beta=1$ y detectar la contradicción
          en la primera entrada del producto: 1 punto.
    \item Concluir que no existe ningún par real $(\alpha,\beta)$ válido: 1 punto.
\end{itemize}
Si $B$ fuera la inversa de $A$, necesariamente $AB=I_2$. Calculamos
\[
    AB=\begin{pmatrix}\alpha&1\\\beta&1\end{pmatrix}
    \begin{pmatrix}\alpha&-1\\-\beta&2\end{pmatrix}
    =\begin{pmatrix}\alpha^2-\beta&2-\alpha\\
        \alpha\beta-\beta&2-\beta\end{pmatrix}.
\]
La segunda columna debe coincidir con $\begin{pmatrix}0\\1\end{pmatrix}$.
Esto obliga a
\[
    2-\alpha=0,\qquad 2-\beta=1,
    \qquad\text{es decir,}\quad \alpha=2,\quad\beta=1.
\]
Pero con estos valores, la primera entrada del producto es
$\alpha^2-\beta=4-1=3\ne1$. Por tanto, $AB\ne I_2$ incluso para el
único par permitido por la segunda columna.
\textbf{No existe ningún par real $(\alpha,\beta)$ para el cual $B=A^{-1}$.}

\section*{4. Inversa por Gauss--Jordan y solución del sistema (6 puntos)}
\noindent\textbf{Rúbrica.}
\begin{itemize}[itemsep=2pt]
    \item Construir $[B\mid I_3]$: 0,5 puntos.
    \item Escalonar con operaciones por filas identificadas: 1,5 puntos.
    \item Reducir hasta $[I_3\mid B^{-1}]$ identificando las operaciones: 2 puntos.
    \item Extraer correctamente $B^{-1}$: 0,5 puntos.
    \item Plantear y calcular $B^{-1}b$: 1 punto.
    \item Presentar y comprobar $(x,y,z)$: 0,5 puntos.
\end{itemize}
\noindent\textbf{Escalonamiento.} Aplicamos las operaciones a toda la fila:
\[
    \left[\begin{array}{ccc|ccc}
        1&0&2&1&0&0\\2&-1&3&0&1&0\\4&1&8&0&0&1
    \end{array}\right]
    \xrightarrow{\substack{F_2\leftarrow F_2-2F_1\\F_3\leftarrow F_3-4F_1}}
    \left[\begin{array}{ccc|ccc}
        1&0&2&1&0&0\\0&-1&-1&-2&1&0\\0&1&0&-4&0&1
    \end{array}\right].
\]
Normalizamos el segundo pivote y anulamos la entrada debajo:
\[
    \xrightarrow{F_2\leftarrow-F_2}
    \left[\begin{array}{ccc|ccc}
        1&0&2&1&0&0\\0&1&1&2&-1&0\\0&1&0&-4&0&1
    \end{array}\right]
    \xrightarrow{F_3\leftarrow F_3-F_2}
    \left[\begin{array}{ccc|ccc}
        1&0&2&1&0&0\\0&1&1&2&-1&0\\0&0&-1&-6&1&1
    \end{array}\right].
\]
El bloque izquierdo está escalonado y tiene tres pivotes.

\noindent\textbf{Reducción.} Normalizamos el tercer pivote y anulamos
las entradas que tiene encima:
\[
    \xrightarrow{F_3\leftarrow-F_3}
    \left[\begin{array}{ccc|ccc}
        1&0&2&1&0&0\\0&1&1&2&-1&0\\0&0&1&6&-1&-1
    \end{array}\right]
    \xrightarrow{\substack{F_1\leftarrow F_1-2F_3\\F_2\leftarrow F_2-F_3}}
    \left[\begin{array}{ccc|ccc}
        1&0&0&-11&2&2\\0&1&0&-4&0&1\\0&0&1&6&-1&-1
    \end{array}\right].
\]
Al obtener $[I_3\mid B^{-1}]$, concluimos que
\[
    B^{-1}=\begin{pmatrix}-11&2&2\\-4&0&1\\6&-1&-1\end{pmatrix}.
\]
Entonces el sistema se resuelve mediante
\[
    \begin{pmatrix}x\\y\\z\end{pmatrix}
    =B^{-1}\begin{pmatrix}2\\-1\\4\end{pmatrix}
    =\begin{pmatrix}-22-2+8\\-8+4\\12+1-4\end{pmatrix}
    =\begin{pmatrix}-16\\-4\\9\end{pmatrix}.
\]
Así, $\boxed{(x,y,z)=(-16,-4,9)}$. La comprobación da
\[
    -16+2(9)=2,\quad 2(-16)-(-4)+3(9)=-1,\quad
    4(-16)+(-4)+8(9)=4.
\]
\end{document}
